\chapter{Segmentation: Architectures, Datasets, and Training}
\label{cap:segmentation}

The segmentation study began with a comparison of FCN, DeepLabV3, and
LR-ASPP on manually annotated lettuce images. Manual segmentation was
time-consuming, so this first comparison was limited to lettuce rather
than repeated for every crop. Its outcome led to the adoption of LR-ASPP.

SAM~3 was then used to generate training labels for lettuce, tomato, and
cucumber images. The resulting multi-crop models differed in their output
classes: plants and path plus background, four classes including covering
material, or path and background alone. These models were evaluated
against the manual lettuce masks. The crop coverage of the training data
therefore differs from that of the reported evaluation set.

This chapter describes the datasets and training procedure, followed by
the offline comparisons that motivate the configuration selected for the
robot. The performance of the trained networks is considered separately
from the quality of the SAM~3 annotations themselves. The architecture results are reported from the available validation
exports; the comparisons of supervision strategies still require checks
on evaluation provenance. Geometric and closed-loop field experiments
are discussed later in the thesis.

\section{Image acquisition and class definition}
\label{sec:c4_acquisition_classes}

\subsection{Image sources and acquisition conditions}
\label{subsec:c4_image_acquisition}

% TODO: Add image sources, camera specifications, field conditions

\subsection{Semantic classes and labeling targets}
\label{subsec:c4_semantic_classes}

Three distinct labeling configurations were evaluated for SAM~3:
\begin{itemize}
\item \textbf{Four-class configuration:} Segments background, plants, soil, and covering material. This configuration was also used for the lettuce models during architecture selection.
\item \textbf{Two-foreground-class configuration:} Labels plants and path, with background as the third output class. This produces three total semantic categories.
\item \textbf{Path-only configuration:} Binary segmentation with background and path as the two output labels. For robotic integration, the path mask is the relevant output. Segmenting soil alone does not identify a unique navigation corridor or establish that it is obstacle-free. The geometric methods in Chapter~\ref{cap:geometric_reference} use the mask to construct a local image-space reference.
\end{itemize}

\section{Dataset preparation}
\label{sec:c4_datasets}

\subsection{Lettuce dataset with manual annotation}
\label{subsec:c4_lettuce}

The manually annotated lettuce dataset, identified as \texttt{Data},
contains 762 image--mask pairs distributed across eight folders. Its
masks were prepared from the manual annotations stored in
\texttt{annotations.json}. The same manual reference is used for the
initial architecture comparison and for evaluating the models trained
with SAM~3-generated multi-crop labels. Automatically annotated lettuce
images belong to the latter training datasets, not to separate
lettuce-only training configurations.

\subsection{SAM~3-assisted annotation for multi-crop datasets}
\label{subsec:c4_sam_annotation}

The time required for manual labeling motivated the use of SAM~3 to
prepare the multi-crop datasets. Initial use suggested faster mask
generation, but the total annotation effort has not been measured in a
systematic comparison that includes inspection and correction.

The SAM~3-assisted labeling procedure covers lettuce, tomato, and
cucumber images for the multi-crop training configurations. In this
workflow, automatic segmentation is used to prepare labels offline; it
is distinct from the segmentation model used during navigation.

For the lettuce images annotated with two foreground classes, the reported
text prompts are \texttt{plants} and \texttt{central farm path}. The
four-class labeling additionally uses \texttt{plastic mulch film} for
covering material. Background accounts for the remaining label in each
configuration. These reported lettuce prompts do not establish the exact
prompts used for every crop in the multi-crop datasets. The binary
configuration retains only the path annotation; crop-specific prompts
and mask-composition settings for each configuration must be documented
from the corresponding labeling runs.

\subsubsection{Tomato dataset}
\label{subsubsec:c4_tomato}

The tomato images were labeled using SAM~3 rather than the manual
procedure used for the initial lettuce comparison. They contribute to
the automatically annotated multi-crop training data. LR-ASPP was adopted
after the comparison of the three candidate architectures on lettuce;
no separate FCN, DeepLabV3, and LR-ASPP comparison was performed on tomato.

\subsubsection{Cucumber dataset}
\label{subsubsec:c4_cucumber}

SAM~3 was also used to label the cucumber images, which contribute to
the automatically annotated multi-crop training data. As for tomatoes,
the architecture choice followed the comparison on manually annotated
lettuce, rather than a separate comparison on cucumber.

\subsection{Dataset partitioning}
\label{subsec:c4_dataset_splits}

The reported training runs use an 80\% training, 10\% validation, and
10\% test partition with seed 42. The available \texttt{train\_model.py}
implements these proportions by assigning the integer parts of the
training and validation counts and leaving the remaining samples for
test. Applied to 762 pairs, this rule gives 609 training, 76 validation,
and 77 test samples. These counts describe the intended manual-dataset
partition; the saved indices and image lists must establish the actual
membership of each evaluated split.

The test set is not used for gradient-based parameter updates. Checkpoint
selection uses validation rather than test performance. A common seed
alone does not ensure that differently sized or ordered datasets yield
identical image partitions. Corresponding image identities must therefore
be checked when comparing the manually supervised lettuce model with the
automatically supervised multi-crop models, including whether a lettuce
evaluation image appears in an automatic model's training set.

The reported operational evaluation selects 42 images from the manual
test split belonging to \texttt{insalata\_grande1},
\texttt{insalata\_grande2}, and \texttt{insalata\_media1}. These folders
represent the row conditions selected for the robot application. The
42 images are a subset of the reported 77 test images, not an additional
independent test set. Performance on this subset describes those selected
conditions, rather than the full range of lettuce images.

\section{Architecture comparison: FCN, DeepLabV3, and LR-ASPP}
\label{sec:c4_architecture_comparison}

This comparison addresses the architecture choice before changing the
annotation procedure or crop coverage. FCN, DeepLabV3, and LR-ASPP were
trained and evaluated using the manually annotated lettuce dataset.

\subsection{Model selection and pretrained weights}
\label{subsec:c4_model_selection}

FCN, DeepLabV3, and LR-ASPP were chosen from the implementations available
in TorchVision. All three were trained on the manually annotated lettuce
dataset. Their selection was practical, rather than the outcome of a
comparison with every available architecture.

The implementation separates each network into a backbone, which extracts
image features, and a segmentation head, which produces class scores at
each output location. Transfer learning is implemented by loading
pretrained parameters and adapting the output layers to the target
labeling task, rather than initializing every network parameter from
scratch.

For FCN and DeepLabV3, the selected implementations use a ResNet-50
backbone with pretrained segmentation weights. After these weights
are loaded, the final convolution of the main classifier is replaced by
a new $1\times1$ convolution with one output channel per target class.

LR-ASPP is initialized differently. The implementation uses a
MobileNetV3-Large backbone and requests its default pretrained weights
through the \texttt{weights\_\allowbreak backbone} argument. However, the
\texttt{weights} argument of the complete segmentation model is set to
\texttt{None}. Consequently, the backbone is pretrained, whereas the
LR-ASPP segmentation head is newly initialized for the requested number
of classes.

This distinction is relevant to the interpretation of the model
comparison. The three configurations differ not only in architecture,
but also in the extent of their pretrained initialization: FCN and
DeepLabV3 retain pretrained segmentation-head layers, whereas LR-ASPP
starts with a new segmentation head.

\subsection{Training configuration}
\label{subsec:c4_training_config}

Training is organized into two stages. Before the first epoch, the
backbone parameters are excluded from gradient-based updates by setting
their \texttt{requires\_grad} attribute to \texttt{False}. The complete
main classifier remains trainable, as does the auxiliary classifier when
present.

For the first ten epochs, AdamW is constructed using only the trainable
parameters, with an initial learning rate of $10^{-3}$ and weight decay
of $10^{-4}$. A cosine-annealing scheduler is stepped after each epoch,
with a minimum learning rate of $10^{-6}$.

At the end of epoch ten, the script enables gradients for all backbone
parameters. Fine-tuning of the backbone therefore begins at epoch eleven,
together with continued training of the segmentation heads. A new AdamW
optimizer is constructed with separate parameter groups: the backbone
starts this stage at a learning rate of $10^{-5}$, while the main and,
where present, auxiliary classifiers start at $10^{-4}$. Weight decay
remains $10^{-4}$.

A checkpoint is saved as \texttt{best.pth} whenever the validation loss is lower than all previous
values; \texttt{last.pth} stores the latest training state. Selection
therefore uses minimum validation loss, not maximum test IoU or simply
the last epoch.

\subsection{Evaluation metrics and results}
\label{subsec:c4_architecture_results}

The comparison uses the three CSV records stored in
\texttt{training insalata}: the FCN, DeepLabV3, and LR-ASPP
\texttt{*\_\allowbreak validation\_\allowbreak metrics.csv} exports.
Each contains aggregate metrics, the checkpoint epoch, and IoU values
for background, plants, soil, and covering material. The reported mIoU
agrees with the arithmetic mean of these four class IoUs, including
background.

Table~\ref{tab:c4_architecture_results} reports the exported metrics.
Inference times have been converted from seconds to milliseconds.
The FPS values equal the reciprocal of the mean recorded time; they
are not measurements of the complete robot processing or control rate.
Checkpoint epochs identify the evaluated weights, not training duration.

\begin{table}[htbp]
    \centering
    \small
    \caption{Architecture comparison from the manual lettuce validation
    exports. Timing values are those recorded in the CSV files, not
    end-to-end robot latency.}
    \label{tab:c4_architecture_results}
    \begin{tabular}{@{}l r r r r r@{}}
        \toprule
        Model & Epoch & mIoU & Path IoU &
        \shortstack{Mean time\\(ms)} & FPS \\
        \midrule
        FCN & 38 & 0.9232 & 0.9621 & 58.24 & 17.17 \\
        DeepLabV3 & 43 & 0.9228 & 0.9651 & 86.87 & 11.51 \\
        LR-ASPP & 194 & 0.9108 & 0.9546 & 10.41 & 96.07 \\
        \bottomrule
    \end{tabular}
\end{table}

FCN and DeepLabV3 achieve mIoU values of 0.9232 and 0.9228, respectively,
compared with 0.9108 for LR-ASPP. Using the unrounded CSV values, the
LR-ASPP deficit is 1.24 percentage points relative to FCN and 1.19
relative to DeepLabV3. Its mean recorded time is 10.41~ms, compared
with 58.24~ms and 86.87~ms. These timings correspond to a 5.60-fold
reduction relative to FCN and an 8.35-fold reduction relative to
DeepLabV3.

Table~\ref{tab:c4_architecture_class_iou} provides the class-level
comparison. LR-ASPP has lower IoU than both alternatives in each class,
although its soil/path IoU remains 0.9546. The data support a trade-off
between accuracy and latency, not a claim of equal segmentation accuracy.

\begin{table}[htbp]
    \centering
    \small
    \caption{Per-class IoU from the same manual lettuce validation
    exports. Plants, soil, and covering material correspond to
    \texttt{pianta}, \texttt{suolo}, and \texttt{copertura} in these files.}
    \label{tab:c4_architecture_class_iou}
    \begin{tabular}{@{}l r r r r@{}}
        \toprule
        Model & Background & Plants & Soil/path & Covering material \\
        \midrule
        FCN & 0.9464 & 0.9280 & 0.9621 & 0.8565 \\
        DeepLabV3 & 0.9464 & 0.9277 & 0.9651 & 0.8519 \\
        LR-ASPP & 0.9394 & 0.9120 & 0.9546 & 0.8372 \\
        \bottomrule
    \end{tabular}
\end{table}

\subsection{Selection of LR-ASPP}
\label{subsec:c4_lraspp_selection}

LR-ASPP was adopted because the substantial reduction in recorded
inference time justified accepting an mIoU approximately 1.2 percentage
points below the other two configurations. It was not the most accurate
model, but offered the practical accuracy--latency compromise sought for
robot perception. FCN and DeepLabV3 had higher IoU but substantially
longer recorded inference times. Subsequent SAM-supervised multi-crop
training therefore used LR-ASPP while varying the output label set.

This choice is based on the available offline comparison. The records
do not establish a hardware-independent speed ranking or demonstrate
that the integrated robot meets its control timing requirements.
Those requirements must be assessed using the deployed model and the
complete perception and control pipeline.

\section{Supervision strategy comparison}
\label{sec:c4_supervision_comparison}

\subsection{Training configurations and evaluation targets}
\label{subsec:c4_supervision_configurations}

Table~\ref{tab:c4_supervision_configurations} summarizes the four training
configurations considered in this comparison. This study is separate
from the initial comparison of the three candidate architectures.
The manual configuration is trained on lettuce only, whereas all three
SAM-supervised configurations use multi-crop data and differ in their
output label sets.

\begin{table}[htbp]
    \centering
    \small
    \caption{Training scope and target labels for the comparison of manual
    and SAM~3-generated supervision. Class counts include background;
    evaluation uses the manual lettuce reference.}
    \label{tab:c4_supervision_configurations}
    \begin{tabularx}{\textwidth}{@{}X X X@{}}
        \toprule
        Configuration & Training scope & Output labels \\
        \midrule
        Manual & Lettuce only (\texttt{Data}) &
        Background, plants, path, covering material \\
        SAM-generated, two foreground classes &
        Multi-crop & Background, plants, path \\
        SAM-generated, four classes &
        Multi-crop & Background, plants, path, covering material \\
        SAM-generated, binary (\texttt{allPlantsSAM}) &
        Multi-crop & Background, path \\
        \bottomrule
    \end{tabularx}
\end{table}

Manual masks supply the ground truth for all comparisons. For path-only
scoring, the target is the manual soil class and the predicted positive
region is the corresponding path output of each model.

\subsection{Evaluation results}
\label{subsec:c4_supervision_results}

\emph{Results to complete after provenance checks:} Present the four-class,
two-foreground-class, and path-only comparisons on the documented lettuce
evaluation set. Report per-class IoU and state which classes each average
includes. Add a path-focused summary, identifying the checkpoint and split
for each configuration.

\subsection{Selection of binary path configuration}
\label{subsec:c4_binary_selection}

The binary configuration was selected because its target is directly
aligned with path extraction and its observed behavior on the selected
operational images supported that choice. The quantitative evidence and
limitations are to be completed in this section. This is not a claim that the
binary configuration outperforms the manually supervised model or that
performance on the selected rows holds across the entire dataset.

\section{Model export and inference configuration}
\label{sec:c4_export_inference}

\subsection{ONNX export}
\label{subsec:c4_onnx_export}

% TODO: Add ONNX export procedure

\subsection{Inference inputs and outputs}
\label{subsec:c4_inference_interface}

% TODO: Add inference interface details
